\section{Cronograma proposto}\label{apx:cronograma}

As apresentações remotas ocorrem em 02 e 04/12~\cite{enunciado}. O enunciado não informa a data de entrega do
relatório, nem em qual dos dois dias cada equipe apresenta; o plano abaixo é uma \textbf{proposta}, que supõe a
entrega do relatório e dos \emph{slides} até a véspera da primeira apresentação (01/12), e deve ser confirmada com
o professor.

\begin{table}[htbp]
\centering
\caption{Cronograma proposto (datas de 2026).}
\label{tab:cronograma}
\small
\begin{tabularx}{\textwidth}{@{}c>{\raggedright\arraybackslash}p{2.3cm}>{\hsize=1.5\hsize}L>{\hsize=.5\hsize}L@{}}
\toprule
 & Período & Atividades & Entrega\\
\midrule
1 & 06--13/11 &
Confirmar com o professor o uso da rede de Hopfield como modelo do tema, as linhas de base e as datas;
dividir as tarefas; baixar o SMD (registrar o \emph{commit}); implementar leitura, divisões e codificação em
termômetro; implementar Hebb, Storkey, recuperação e escores; executar E1a. &
Dados e rede clássica validados em E1a\\
2 & 16--20/11 (20/11: feriado) &
Seleção de protótipos e E1b; detector nas máquinas 1-3, 2-5 e 3-6; linhas de base (\emph{Random Forest}, SVM,
OC-SVM e \emph{k}-NN); escolha do limiar em $L_{\mathrm{val}}$. &
Primeiros resultados da Etapa~I nas 3 máquinas\\
3 & 23--27/11 &
Protocolo~S (18 máquinas) e Protocolo~U (28); ablações; rede moderna ($\beta$); Etapa~II (\emph{zero-shot} e
\emph{few-shot}, com registro de tempo e custo); preencher a Seção~\ref{sec:resultados}. &
Seção~\ref{sec:resultados} completa\\
4 & 30/11--01/12 &
Revisar as Seções~\ref{sec:critica} e~\ref{sec:conclusoes} e o resumo à luz dos resultados; conferir as
referências; preparar os \emph{slides} e ensaiar. &
Relatório e \emph{slides}\\
5 & 02 e 04/12 & Apresentação remota. & ---\\
\bottomrule
\end{tabularx}
\end{table}

\paragraph{Divisão de tarefas (equipe de até 3 integrantes).}
\begin{itemize}
  \item \textbf{Integrante A} (\apreencher[nome]): dados, divisões, linhas de base clássicas e métricas.
  \item \textbf{Integrante B} (\apreencher[nome]): rede de Hopfield (clássica e moderna), E1 e ablações.
  \item \textbf{Integrante C} (\apreencher[nome]): Etapa~II (\emph{prompt}, execução e custo) e revisão do texto.
\end{itemize}

\paragraph{Riscos e mitigação.}
\begin{itemize}
  \item \textbf{Hopfield próximo do acaso:} o resultado é registrado e analisado como achado (ablações e
        Seção~\ref{sec:critica}); a comparação crítica vale igualmente.
  \item \textbf{Prazo apertado:} priorizar o Protocolo~S nas 3 máquinas representativas e estender às 18 depois.
  \item \textbf{Custo ou latência da API:} reduzir $n$, usar o modelo mais barato disponível e limitar $k=20$ a
        uma semente.
  \item \textbf{Modelo fora da lista do enunciado:} a confirmação com o professor é a primeira atividade do
        cronograma.
\end{itemize}
